% !TeX program = xelatex

\documentclass[aspectratio=169]{beamer}

% ============================================================
% 1. Fonts
% ============================================================
\usepackage{fontspec}

\setmainfont{Microsoft JhengHei}
\setsansfont{Microsoft JhengHei}
\setmonofont{JetBrains Mono}

\XeTeXlinebreaklocale "zh"


% ============================================================
% 2. Mathematics, enumerate
% ============================================================
\usepackage{amsmath}
\usepackage{amssymb}
\usepackage{enumitem}

\setlistdepth{6}
\renewlist{itemize}{itemize}{9}
\renewlist{enumerate}{enumerate}{9}
\setlist[itemize,1]{label=\textbullet}
\setlist[itemize,2]{label=$\circ$}
\setlist[itemize,3]{label=$\ast$}
\setlist[itemize,4]{label=$\cdot$}
\setlist[itemize,5]{label=--}
\setlist[itemize,6]{label=--}
\setlist[enumerate,1]{label=\arabic*.}
\setlist[enumerate,2]{label=\arabic*.}
\setlist[enumerate,3]{label=\arabic*.}
\setlist[enumerate,4]{label=\arabic*.}


% ============================================================
% 3. Images / Tables / Boxes
% ============================================================
\usepackage{graphicx}
\usepackage{booktabs}
\usepackage[most]{tcolorbox}


% ============================================================
% 4. Colors
% ============================================================
\definecolor{myRed}{RGB}{145,90,90}
\definecolor{myGreen}{RGB}{85,115,100}
\definecolor{myBlue}{RGB}{70,100,130}
\definecolor{myOrange}{RGB}{155,115,80}
\definecolor{myPurple}{RGB}{115,100,130}

\definecolor{myGray}{RGB}{242,242,238}
\definecolor{myDarkgray}{RGB}{90,90,88}
\definecolor{myBlack}{RGB}{45,45,43}


% ============================================================
% 5. Beamer Colors
% ============================================================
\setbeamercolor{structure}   {fg=myBlue}
\setbeamercolor{title}       {fg=white, bg=myBlue}
\setbeamercolor{frametitle}  {fg=white, bg=myBlue}
\setbeamercolor{block title} {fg=white, bg=myBlue}
\setbeamercolor{block body}  {fg=myBlack, bg=myGray}
\setbeamercolor{alerted text}{fg=myRed}
\setbeamercolor{example text}{fg=myGreen}


% ============================================================
% 6. Theme
% ============================================================
\usetheme{Madrid}
\usecolortheme{default}

\setbeamertemplate{navigation symbols}{}


% ============================================================
% 7. Footer
% ============================================================
\setbeamertemplate{footline}{
    \leavevmode
    \hbox{
        % Left: Author
        \begin{beamercolorbox}[
            wd=.33\paperwidth,
            ht=2.5ex,
            dp=1ex,
            center
        ]{author in head/foot}
            \insertdate
        \end{beamercolorbox}

        % Center: Title
        \begin{beamercolorbox}[
            wd=.33\paperwidth,
            ht=2.5ex,
            dp=1ex,
            center
        ]{title in head/foot}
            \insertshorttitle
        \end{beamercolorbox}

        % Right: Page Number
        \begin{beamercolorbox}[
            wd=.33\paperwidth,
            ht=2.5ex,
            dp=1ex,
            center
        ]{date in head/foot}
            \insertframenumber{} / \inserttotalframenumber
        \end{beamercolorbox}
    }
}


% ============================================================
% 8. Presentation Information
% ============================================================
\title[第二次雙周會]{畢業專題 文獻探討}
\subtitle{Agentic AI Security}
\author{Syu}
\date{October 5, 2026}



% ============================================================
% Helpers (table spacing)
% ============================================================
\newcommand{\tabsetup}{\renewcommand{\arraystretch}{1.18}}


% ========================= DocStart =========================
\begin{document}

\begin{frame}
    \titlepage
\end{frame}

\begin{frame}{Index}
    \tableofcontents
\end{frame}

\section{回顧}
\begin{frame}{回顧}
    \begin{enumerate}
        \item A Defensive Framework Against Adversarial Attacks on Machine Learning-Based Network Intrusion Detection Systems (arXiv 2502.15561)
        \item Adversarial Challenges in Network Intrusion Detection Systems: Research Insights and Future Prospects (arXiv 2409.18736)
    \end{enumerate}
\end{frame}


% ############################################################
% Paper 1 : Agentic AI Security (Datta et al.)
% ############################################################
\section{Agentic AI Security (arXiv 2510.23883)}

\begin{frame}{Agentic AI Security (arXiv 2510.23883)}
    \small
    \textbf{Agentic AI Security: Threats, Defenses, Evaluation, and Open Challenges}\\
    Chhabra, Datta, Nahin, Mohapatra（University of South Florida）

    \medskip

    \begin{block}{摘要}
        Agentic AI 具備規劃、工具使用、記憶與自主性，能在網頁、軟體與實體環境中自動執行任務，
        也因此產生不同於傳統 AI 安全與軟體安全的新風險。本文為 survey：整理\alert{威脅分類}、\alert{防禦策略}、\alert{評測方法}與\alert{開放問題}，目標是支援 secure-by-design 的 agent 系統。
    \end{block}

    \begin{block}{論文結構}
        威脅分類 (§3) $\rightarrow$ 防禦與安全控制 (§4) $\rightarrow$ Benchmark 與評估 (§5) $\rightarrow$ 開放問題 (§6)
    \end{block}
\end{frame}


\begin{frame}{動機：為什麼 Agentic AI 的安全不同}
    \small
    \begin{itemize}
        \item 一般 LLM 多為被動回應；Agent 具備持久記憶、跨時間推理、與其他 agent 協作
        \item 自主性與持久性 $\rightarrow$ 攻擊面擴大；工具整合 $\rightarrow$ 放大濫用；多 agent 協作 $\rightarrow$ 行為難預測
    \end{itemize}

    \begin{block}{真實案例}
        \begin{itemize}
            \item \textbf{EchoLeak (CVE-2025-32711)}：惡意郵件中的 prompt 讓 Microsoft Copilot 在無使用者互動下外洩資料
            \item Symantec 以 OpenAI Operator 實驗：可蒐集個資並自動化 credential stuffing
        \end{itemize}
    \end{block}

    \begin{alertblock}{Lupinacci et al. 的測試結果}
        94.4\% 的 SOTA LLM agent 易受 prompt injection，83.3\% 易受 retrieval-based backdoor，100\% 易受 inter-agent trust exploit
    \end{alertblock}
\end{frame}


\begin{frame}{威脅分類總覽}
    \tabsetup
    \footnotesize
    \begin{table}
        \centering
        \begin{tabular}{@{}l p{8.2cm}@{}}
            \toprule
            類別 & 主要內容 \\
            \midrule
            Prompt Injection \& Jailbreaks & 直接/間接、有意/無意、依模態、依傳播性、多語言與混淆、payload splitting \\
            Autonomous Cyber-Exploitation \& Tool Abuse & one-day 漏洞利用、自主網站入侵、emergent tool abuse \\
            Multi-agent \& Protocol-level & MCP / A2A 協定攻擊、冒充、協作操弄、責任歸屬模糊 \\
            Interface \& Environment Risks & 觀察/行動空間錯位、perception-action 脆弱性、動態內容與 CAPTCHA \\
            Governance \& Autonomy & 人類監督不足、自主等級、治理框架 \\
            \bottomrule
        \end{tabular}
    \end{table}
    \small 各類別並非互斥，依主要特性歸類以利分析。
\end{frame}


\begin{frame}{Prompt Injection：直接與間接}
    \footnotesize
    \begin{columns}[T]
        \begin{column}{0.48\textwidth}
            \begin{block}{Direct (DPI)}
                \begin{itemize}
                    \item 惡意指令直接放入 prompt
                    \item 攻擊者 = 終端使用者
                    \item 例：「ignore all previous instructions…」，誘使客服 bot 洩漏內部資料或寄信
                \end{itemize}
            \end{block}
        \end{column}
        \begin{column}{0.48\textwidth}
            \begin{block}{Indirect (IPI)}
                \begin{itemize}
                    \item 惡意指令藏在網頁、API、資料庫等外部資料
                    \item 攻擊者 = 第三方資料提供者
                    \item 更隱蔽：看起來像合法指令
                    \item Adaptive attack 對 8 種 IPI 防禦的成功率約 50\%
                \end{itemize}
            \end{block}
        \end{column}
    \end{columns}

    \medskip

    \tabsetup
    \begin{tabular}{@{}l p{9.3cm}@{}}
        \toprule
        其他分類維度 & 說明 \\
        \midrule
        有意 / 無意 & 含糊的使用者輸入或長對話的 context drift 也可能覆蓋系統指令 \\
        模態 & 文字（code / P2SQL）、圖像/影片、音訊、混合（CrossInject 效果至少提升 30.1\%） \\
        傳播性 & 非傳播 vs. 傳播（recursive injection、AI worm） \\
        混淆與拆分 & Base64、emoji、低資源語言；payload splitting 將惡意內容拆成多段，聚合時才成形 \\
        \bottomrule
    \end{tabular}
\end{frame}


\begin{frame}{自主網路利用與工具濫用}
    \small
    \begin{itemize}
        \item \textbf{One-day 漏洞}：GPT-4 agent 在提供 CVE 描述時成功率 87\%，優於 OWASP ZAP、Metasploit 等傳統工具
        \item \textbf{自主網站入侵}：在不預知漏洞下完成 XSS+CSRF 串接、SSTI、blind SQL union injection；缺少文件存取或詳細 prompt 時效能大降
        \item \textbf{Emergent tool abuse}：ConAgents（多 agent 分工選工具、執行、校正）、ToolGen（把工具呼叫放進 next-token 生成）使工具串接更容易
    \end{itemize}

    \begin{alertblock}{攻擊經濟學}
        每次 one-day exploit 成本僅數美元，且可平行化，大規模攻擊比雇用人類攻擊者便宜
    \end{alertblock}

    $\Rightarrow$ 需要 containment 與 runtime monitoring
\end{frame}


\begin{frame}{多代理與協定層威脅}
    \footnotesize
    \begin{columns}[T]
        \begin{column}{0.48\textwidth}
            \begin{block}{MCP（client--server）}
                \begin{itemize}
                    \item Flooding / replay $\rightarrow$ DoS
                    \item Token 外洩或不安全 proxy $\rightarrow$ 冒充 agent
                    \item Embedded backdoor、retrieval corruption
                    \item 推論與 side-channel 導致機密外洩
                \end{itemize}
            \end{block}
        \end{column}
        \begin{column}{0.48\textwidth}
            \begin{block}{A2A（agent--agent）}
                \begin{itemize}
                    \item 假 agent 宣告、未授權註冊
                    \item 遞迴任務委派 $\rightarrow$ deadlock / 無限迴圈 DoS
                    \item Transitive prompt injection
                    \item Context tampering、memory manipulation
                \end{itemize}
            \end{block}
        \end{column}
    \end{columns}

    \medskip

    \begin{block}{Threat actor 視角的六類威脅（改編自 Ko et al.）}
        Impersonation \& Role Abuse、Coordination Manipulation、Knowledge \& Learning Manipulation、Inference \& Policy Evasion、Accountability Obfuscation、Confidentiality \& Integrity Breaches
    \end{block}
\end{frame}


\begin{frame}{介面環境風險與治理}
    \small
    \begin{itemize}
        \item \textbf{觀察/行動空間錯位}：LLM 以靜態文字預訓練，卻要做 scroll、hover、切分頁；AgentOccam 簡化動作空間後穩定度提升
        \item \textbf{Perception-action 脆弱}（WebArena）：
        \begin{itemize}
            \item 誤解先前輸入、重複相同搜尋直到步數用盡
            \item 提前終止：GPT-4 將 54.9\% 可完成的任務誤判為不可能
            \item 模板脆弱：僅在 61 個模板中的 4 個穩定成功
        \end{itemize}
        \item \textbf{動態內容與 CAPTCHA}：多模態 agent 最高約 40\% vs. 人類近 100\%
        \item \textbf{治理與自主性}：低監督下可能出現不可預測行為與劫持；建議 human-in-the-loop 與結構化自主等級
    \end{itemize}
\end{frame}


\begin{frame}{防禦 (1)：Prompt Injection 抵抗設計}
    \tabsetup
    \footnotesize
    \begin{table}
        \centering
        \begin{tabular}{@{}l p{4.6cm} p{4.6cm}@{}}
            \toprule
            類型 & 代表方法 & 取捨 \\
            \midrule
            Agent-focused & Instruction hierarchy、StruQ、SecAlign (DPO)、circuit breaking & 需重訓，可能降低通用能力，對 adaptive attack 效果有限 \\
            System-focused & Dual-LLM、Plan-then-Execute、Melon、TrustAgent；偵測、隔離、prompt augmentation、perplexity & 保留模型通用性但系統更複雜，多階段/多 agent 攻擊仍可繞過 \\
            User-focused & 人類驗證 & 無法擴展，攻擊者就是使用者時失效 \\
            \bottomrule
        \end{tabular}
    \end{table}

    \begin{block}{結論}
        完全阻止 prompt injection 仍是開放問題；實務上傾向\alert{混合式 (hybrid) 設計}，組合不同類型的防禦。
    \end{block}
\end{frame}


\begin{frame}{防禦 (2)：政策執行、沙箱、監控、標準}
    \footnotesize
    \begin{itemize}
        \item \textbf{Policy enforcement}
        \begin{itemize}
            \item Governance-centric（runtime）：GuardAgent、AgentSpec（DSL 定義規則）、ShieldAgent、R$^2$-Guard
            \item Signal-centric（掃描輸入輸出）：Llama Guard、LlavaGuard、Safewatch
        \end{itemize}
        \item \textbf{Sandboxing}：SandboxEval、容器化（gVisor / Docker）、ToolEmu；限制：沙箱本身有漏洞（跨 session 洩漏）、降低效率
        \item \textbf{Detection \& monitoring}：Pro2Guard 以機率可達性分析預測違規；Adversarial Markov Games 讓攻防共同演化
        \item \textbf{標準}：NIST AI RMF GenAI Profile、OWASP Agentic AI Threats、CSA MAESTRO
    \end{itemize}

    \begin{alertblock}{Table 1 的觀察}
        21 個方法中，涵蓋 Rate limiting 與 Rollback/Recovery 的各只有 3 個，多數集中在輸入治理與政策層。
    \end{alertblock}
\end{frame}


\begin{frame}{Benchmark 與評估}
    \footnotesize
    \begin{columns}[T]
        \begin{column}{0.48\textwidth}
            \begin{block}{Benchmark 類型}
                \begin{itemize}
                    \item 能力導向：BrowserGym、$\tau$-bench、OSWorld、AndroidWorld
                    \item 安全導向：ST-WebAgentBench、AgentHarm、OS-Harm、R-Judge、ToolEmu、AgentDojo、InjecAgent、ASB、\alert{Agent-SafetyBench}
                \end{itemize}
            \end{block}
        \end{column}
        \begin{column}{0.48\textwidth}
            \begin{block}{評估方式的演進}
                \begin{itemize}
                    \item Process-aware：評分軌跡而非只看最終結果
                    \item 重複試驗：pass$^{\wedge k}$（GPT-4o 在 retail 的 pass$^{\wedge 8}$ 可低於 25\%）
                    \item 標準化 LLM judge：回報與人類一致性
                    \item 沙箱/模擬器需說明擬真度
                    \item 可重現性：釋出攻擊模板、回報成本與延遲
                \end{itemize}
            \end{block}
        \end{column}
    \end{columns}
\end{frame}


\begin{frame}{開放問題與結論}
    \footnotesize
    \begin{columns}[T]
        \begin{column}{0.48\textwidth}
            \begin{enumerate}
                \item \textbf{長時程安全}：時間穩健性、潛伏式不當行為偵測
                \item \textbf{多 agent 安全}：單一惡意 agent 可造成連鎖失敗；需認證通訊、回滾機制、sentinel agent
                \item \textbf{更好的 benchmark}：分布式指標、可靠的 LLM judge、高擬真沙箱
            \end{enumerate}
        \end{column}
        \begin{column}{0.48\textwidth}
            \begin{enumerate}
                \setcounter{enumi}{3}
                \item \textbf{Adaptive attack}：多數防禦只做靜態評估
                \item \textbf{實體世界}：感測器欺騙（GPS、LiDAR）、端到端 VLA 模型缺乏防禦分層
                \item \textbf{人機介面}：人類監督本身可被社交工程影響
            \end{enumerate}
        \end{column}
    \end{columns}
\end{frame}


% ############################################################
% Paper 2 : Agentic AI and Cybersecurity (Lazer et al.)
% ############################################################
\section{Agentic AI and Cybersecurity (arXiv 2601.05293)}

\begin{frame}{Agentic AI and Cybersecurity (arXiv 2601.05293)}
    \small
    \textbf{A Survey of Agentic AI and Cybersecurity: Challenges, Opportunities and Use-case Prototypes}\\
    Lazer, Aryal, Gupta, Bertino（Tennessee Tech / Univ. of Nebraska Omaha / Purdue）

    \medskip

    \begin{block}{問題}
        Agentic AI 在資安上是\alert{雙面刃}：能強化持續監控、事件回應、威脅獵捕；同樣的能力也加速偵察、利用與社交工程。
        現有治理與保證機制多假設短生命週期、有人在迴圈的 AI，無法涵蓋 collusion、cascading failure、memory poisoning、oversight evasion 等系統性風險。
    \end{block}

    \begin{block}{貢獻}
        概念基礎 $\rightarrow$ 防禦應用 $\rightarrow$ 攻擊應用 $\rightarrow$ 系統性風險 $\rightarrow$ 框架與 benchmark $\rightarrow$ 量子考量 $\rightarrow$ 三個原型實作
    \end{block}
\end{frame}


\begin{frame}{什麼是 Agentic AI：架構與自主等級}
    \footnotesize
    \begin{columns}[T]
        \begin{column}{0.42\textwidth}
            \begin{block}{核心特性}
                \begin{itemize}
                    \item Reasoning（CoT、Reflection）
                    \item Interaction（工具、API、程式執行）
                    \item Autonomy
                    \item Adaptability（記憶與回饋）
                \end{itemize}
            \end{block}
            以 foundation model 為核心，加上記憶、檢索 (RAG)、工具，形成規劃--行動--反思的迴圈；多 agent 架構由 coordinator 分派給 task agents。
        \end{column}
        \begin{column}{0.54\textwidth}
            \tabsetup
            \begin{tabular}{@{}c p{2.2cm} p{2.6cm}@{}}
                \toprule
                Level & 自主程度 & 資安意涵 \\
                \midrule
                0 & Static inference & 風險最低、可預測 \\
                1 & Assistive & 低風險 \\
                2 & Tool-assisted (RAG) & 中等風險 \\
                3 & Semi-agentic & 高風險 \\
                4 & Fully agentic & 極高、難稽核 \\
                \bottomrule
            \end{tabular}
        \end{column}
    \end{columns}
\end{frame}


\begin{frame}{防禦應用}
    \tabsetup
    \footnotesize
    \begin{table}
        \centering
        \begin{tabular}{@{}l p{7.3cm}@{}}
            \toprule
            領域 & 內容 \\
            \midrule
            Autonomous cyber defense & 自主監控與偵測、事件回應、協同防禦 \\
            Threat intelligence \& adversarial analysis & 主動威脅獵捕、攻擊模擬（滲透測試、CTF）、對抗式 AI 防禦（釣魚等） \\
            Enterprise automation \& governance & DevSecOps、IAM、資安人力擴增 \\
            Simulation, training \& testing & 自動生成 cyber range（ARCeR）、高擬真訓練環境（Cyberwheel） \\
            \bottomrule
        \end{tabular}
    \end{table}

    \small
    \begin{itemize}
        \item 多為\alert{模組化 multi-agent}，各 agent 對應 NIST 單一功能（Govern / Identify / Protect / Detect / Respond / Recover）
        \item 取捨：速度 vs. 執行權限；不可逆動作需人類授權（Knack \& Burke）
        \item 定位為\alert{擴增而非取代}分析師（全球資安人力缺口約 4--5 百萬）
    \end{itemize}
\end{frame}


\begin{frame}{攻擊應用：Agentic AI 帶來的攻擊}
    \tabsetup
    \footnotesize
    \begin{table}
        \centering
        \begin{tabular}{@{}l p{8.3cm}@{}}
            \toprule
            攻擊領域 & 重點 \\
            \midrule
            Ransomware & Unit 42 框架模擬完整攻擊鏈約 25 分鐘；外洩平均時間由 2021 年 9 天降至 2024 年約 2 天 \\
            Insider threat & 被入侵的 agent 以合法憑證長期持續執行看似正常的動作 \\
            Social engineering \& fraud & ScamAgents 可進行多輪詐騙電話並規避 guardrail \\
            Exploitation \& recon & HexStrike-AI 被濫用於平行掃描數千個 IP \\
            \bottomrule
        \end{tabular}
    \end{table}

    \begin{block}{論文觀點}
        Agentic ransomware 是\alert{自動化與決策自主的升級}，而非新的密碼學或漏洞類別；調整發生在 orchestration 層。重點轉向「在傳播與外洩完成前打斷自主攻擊迴圈」。
    \end{block}
\end{frame}


\begin{frame}{Agentic AI 本身的安全：四層風險模型}
    \tabsetup
    \footnotesize
    \begin{table}
        \centering
        \begin{tabular}{@{}l p{4.4cm} p{4.6cm}@{}}
            \toprule
            層 & 威脅 & 防禦 \\
            \midrule
            Perception & Data poisoning、adversarial input、supply chain & Input validation、guardrails、anomaly detection \\
            Reasoning & Prompt injection、logic manipulation、backdoor & Multi-agent debate、reasoning guardrails \\
            Action & API 濫用、unsafe tool invocation、未授權程式執行 & Sandboxing、privilege separation \\
            Memory & Poisoned retrieval、資料外洩 & 加密、細粒度存取控制 \\
            \bottomrule
        \end{tabular}
    \end{table}

    \small
    來源：Wong \& Saade。失敗會跨層傳播，因此需要\alert{跨層協同}的控制。
    其他框架：ATFAA / SHIELD、MAESTRO、OWASP Agentic Security Initiative、Microsoft 失敗模式分類。
\end{frame}


\begin{frame}{系統層漏洞與 Prompt Injection 緩解}
    \footnotesize
    \begin{block}{He et al.：三類系統層漏洞}
        Session 管理失敗、model pollution 與隱私外洩、可執行 agent 程式導致任意程式碼執行。
        Bash agent 實驗：未使用沙箱時超過 75\% 惡意指令可成功執行，容器沙箱則幾乎全部擋下。
    \end{block}

    \tabsetup
    \begin{tabular}{@{}l p{3.4cm} p{4.4cm}@{}}
        \toprule
        緩解類型 & 機制 & 限制 \\
        \midrule
        Input sanitization & 擋明顯樣式的注入 & 對語意、間接、多模態注入無效 \\
        Least-privilege 憑證 & 縮小爆炸半徑 & 無法阻止在允許範圍內被誤導 \\
        Runtime monitoring / red teaming & 偵測不安全行為 & 多為事後，難抓跨工具連鎖失敗 \\
        Intent-bound delegation & 行動綁定授權意圖 & 依賴上游意圖正確指定 \\
        \bottomrule
    \end{tabular}
\end{frame}


\begin{frame}{多代理、資源濫用、推理與記憶}
    \footnotesize
    \begin{columns}[T]
        \begin{column}{0.48\textwidth}
            \begin{block}{多代理安全}
                \begin{itemize}
                    \item 共享記憶/權限讓單一被入侵 agent 引發連鎖；隱寫術可做隱蔽協作
                    \item 推理層防禦：PeerGuard、debate / voting
                    \item 基礎設施層：GUARDIAN、SentinelAgent（圖形異常偵測）
                    \item BlockA2A、SAFEFLOW 做身分與資訊流控制
                    \item \alert{Multi-agent security tax}：控制愈強，協作效率愈低
                \end{itemize}
            \end{block}
        \end{column}
        \begin{column}{0.48\textwidth}
            \begin{block}{其他風險}
                \begin{itemize}
                    \item \textbf{資源濫用}：CRABS 類攻擊以合法輸入引發過量 token 與工具呼叫，成本放大式 DoS
                    \item \textbf{推理劫持}：竄改 reasoning trace 重導多步規劃
                    \item \textbf{記憶污染}：污染可在攻擊結束後仍影響未來任務
                    \item 人機層面：agent 代為點擊、下載、輸入憑證，擴大社交工程面
                \end{itemize}
            \end{block}
        \end{column}
    \end{columns}
\end{frame}


\begin{frame}{原型實作（Mistral via Ollama，合成資料）}
    \tabsetup
    \footnotesize
    \begin{table}
        \centering
        \begin{tabular}{@{}l p{3.2cm} p{5.4cm}@{}}
            \toprule
            原型 & 設定 & 結果 \\
            \midrule
            Red / Blue agents & 讀系統設定檔，兩輪攻防 & 雙方依對手上一輪行為自行調整策略 \\
            SOC triage agent & 6 筆合成警報、policy gate 限制唯讀 & Baseline：P 1.00, R 0.75, F1 0.86, MTTR 7.5 分；Agent：P / R / F1 皆 1.00, MTTR 6.0 分 \\
            Anomaly / IDS agent & 合成 NetFlow + RAG notes & Baseline（z-score + IForest）：AUROC 0.87, F1 0.74；Agent：AUROC 0.93, F1 0.81 \\
            \bottomrule
        \end{tabular}
    \end{table}

    \small
    \textbf{個人註記}：實驗為可行性展示，資料為合成且規模很小（例如只有 6 筆警報），數字不宜外推到真實環境。
\end{frame}


\begin{frame}{Benchmark、量子與未來方向}
    \footnotesize
    \begin{columns}[T]
        \begin{column}{0.48\textwidth}
            \begin{block}{Benchmark 與量子}
                \begin{itemize}
                    \item 系統級（BountyBench）、情境式（ARCeR、RedTeamLLM）、領域式（FinGAIA）、聚焦式（WASP、AgentVigil）
                    \item 取捨：廣度 vs. 診斷精準度；自動化 vs. 擬真度
                    \item 量子：quantum agents、QML 偵測、後量子密碼（harvest-now-decrypt-later 風險）
                \end{itemize}
            \end{block}
        \end{column}
        \begin{column}{0.48\textwidth}
            \begin{block}{未來方向}
                \begin{itemize}
                    \item 即時輸入檢查（區分指令與資料）
                    \item 內部執行透明度與結構化 log
                    \item 帶形式化約束的規劃
                    \item 與不可信環境的安全互動
                    \item 多 agent 協作安全、記憶完整性
                    \item 統一 benchmark、治理與部署保證
                \end{itemize}
            \end{block}
        \end{column}
    \end{columns}

    \medskip
    \small \textbf{結論}：自主性提升速度與適應力，但降低可預測性與可稽核性；需要有界自主、可強制的執行控制、持續保證與人類監督。
\end{frame}


% ############################################################
% Paper 3 : Agent-SafetyBench (Zhang et al.)
% ############################################################
\section{Agent-SafetyBench (arXiv 2412.14470)}

\begin{frame}{Agent-SafetyBench (arXiv 2412.14470)}
    \small
    \textbf{Agent-SafetyBench: Evaluating the Safety of LLM Agents}\\
    Zhang, Cui, Lu, Zhou et al.（Tsinghua University, CoAI group）

    \medskip

    \begin{block}{問題}
        過去研究多關注 LLM 的\alert{內容安全}（是否產生不安全文字）。當 LLM 作為 agent 與環境、工具互動，出現新的\alert{行為安全}問題，例如在公開論壇洩漏敏感資訊、誤改訂單數量。
        這類行為不一定在文字上有害，需要理解風險才能避免，但缺乏全面的評測。
    \end{block}

    \begin{block}{做法}
        提出 Agent-SafetyBench：349 個互動環境、2,000 個測試案例、8 類風險、10 種失敗模式，並微調一個專用 scorer 自動判斷安全性。
    \end{block}
\end{frame}


\begin{frame}{Benchmark 概覽與資料建構}
    \footnotesize
    \begin{columns}[T]
        \begin{column}{0.48\textwidth}
            \begin{block}{特色}
                \begin{itemize}
                    \item 349 個環境（先前最多 53 個）、每類風險 250 筆，共 2,000 筆
                    \item 大量\alert{沒有公開 API} 的新環境（例：MindCloning、NanorobotController）
                    \item 每個案例標註預期的失敗模式
                    \item 環境為 JSON 工具 schema + Python class，可自訂初始參數
                \end{itemize}
            \end{block}
        \end{column}
        \begin{column}{0.48\textwidth}
            \begin{block}{資料來源}
                \begin{itemize}
                    \item 精修既有資料集（R-Judge、AgentDojo、ToolEmu 等）得 876 筆（含由 AdvBench 與 20 種越獄模板組成的 200 筆）
                    \item GPT-4o 擴增得 1,124 筆：先生成 300 個新環境名稱以提升多樣性，再生成預期風險行為以提升品質
                    \item 品質控管：人工預檢、自動驗證、人工後檢；交叉驗證 98\% 案例、97.5\% 標籤合理
                \end{itemize}
            \end{block}
        \end{column}
    \end{columns}
\end{frame}


\begin{frame}{8 類風險與 10 種失敗模式}
    \footnotesize
    \begin{columns}[T]
        \begin{column}{0.44\textwidth}
            \begin{block}{風險類別}
                \begin{enumerate}
                    \item 洩漏敏感資料
                    \item 造成財產損失
                    \item 散布不安全資訊 / 假訊息
                    \item 造成人身傷害
                    \item 違法 / 違反倫理
                    \item 損害可用性
                    \item 有害 / 有漏洞程式碼
                    \item 產生不安全資訊（不用工具）
                \end{enumerate}
            \end{block}
        \end{column}
        \begin{column}{0.54\textwidth}
            \begin{block}{失敗模式 (M1--M10)}
                \begin{enumerate}
                    \item[M1] 不呼叫工具直接產生有害內容
                    \item[M2] 資訊不完整仍呼叫工具
                    \item[M3] 取得完整資訊前就呼叫工具
                    \item[M4] 忽略已知限制條件
                    \item[M5] 忽略隱含或潛在風險
                    \item[M6] 使用錯誤參數
                    \item[M7] 忽略工具已知問題仍呼叫
                    \item[M8] 未呼叫必要工具
                    \item[M9] 過度信任工具結果
                    \item[M10] 無法從多個結果挑出安全選項
                \end{enumerate}
            \end{block}
        \end{column}
    \end{columns}
\end{frame}


\begin{frame}{評分方式：微調 Scorer}
    \small
    \begin{itemize}
        \item 行為安全需考慮環境中的互動，直接用 GPT-4o 當 judge 在 200 筆 Gemini-1.5-Flash 紀錄上準確率僅 \alert{75.5\%}
        \item 以人工標註的 4,000 筆互動紀錄（2,186 unsafe / 1,814 safe）訓練，GPT-4o 針對人工標籤生成分析說明（抽查 94\% 合理）
        \item 基礎模型 Qwen2.5-7B-Instruct，輸出 safe / unsafe 與分析，準確率 \alert{91.5\%}
        \item 評測流程：agent 選工具 $\rightarrow$ 環境回傳結果 $\rightarrow$ 重複，直到給出最終回覆；安全分數 = 判為 safe 的比例
    \end{itemize}
\end{frame}


\begin{frame}{主要結果：16 個 agent 皆低於 60\%}
    \tabsetup
    \footnotesize
    \begin{columns}[T]
        \begin{column}{0.50\textwidth}
            \begin{tabular}{@{}l r r r@{}}
                \toprule
                Model & Total & Behavior & Content \\
                \midrule
                Claude-3-Opus & 59.8 & 53.2 & 84.9 \\
                Claude-3.5-Sonnet & 59.4 & 51.9 & 88.6 \\
                Claude-3.5-Haiku & 55.1 & 40.7 & 86.4 \\
                GPT-4o & 44.2 & 36.9 & 72.5 \\
                Gemini-1.5-Flash & 41.6 & 34.6 & 69.1 \\
                Llama3.1-405B & 35.4 & 24.0 & 79.6 \\
                Qwen2.5-7B & 18.8 & 13.5 & 38.9 \\
                \midrule
                Average (16) & 38.5 & 30.4 & 68.4 \\
                \bottomrule
            \end{tabular}
        \end{column}
        \begin{column}{0.46\textwidth}
            \begin{itemize}
                \item 強模型整體較安全；閉源優於開源
                \item 行為安全明顯弱於內容安全（即使多數行為案例沒有越獄攻擊）
                \item 各類風險平均：Spread 僅 15.6\%（agent 常不驗證就透過貼文、郵件散布資訊）；Produce 87.0\% 最高（越獄已被大量研究）
            \end{itemize}
        \end{column}
    \end{columns}
\end{frame}


\begin{frame}{失敗模式與有用性分析}
    \footnotesize
    \begin{columns}[T]
        \begin{column}{0.48\textwidth}
            \begin{block}{各失敗模式平均安全分數}
                \begin{itemize}
                    \item 較高：M1 70.1、M6 63.3、M10 63.1
                    \item 最低：M7 12.5、M2 18.1（捏造參數、忽略工具危險）
                    \item 偏低：M5 23.2、M3 27.0、M4 28.3
                    \item M9 37.5：只有單一選項時缺乏驗證
                \end{itemize}
            \end{block}
        \end{column}
        \begin{column}{0.48\textwidth}
            \begin{alertblock}{兩個根本缺陷}
                \begin{enumerate}
                    \item \textbf{缺乏穩健性}：工具使用不可靠（例如數量填錯）
                    \item \textbf{缺乏風險意識}：參數正確卻忽略後果（例如關閉所有警報）
                \end{enumerate}
            \end{alertblock}
        \end{column}
    \end{columns}

    \medskip

    \small
    將案例分為可安全完成 (fulfillable) 與不可 (unfulfillable)：強安全模型在 fulfillable 案例仍保持相當有用性（不是靠拒絕），在 unfulfillable 案例則更少提供協助，顯示較強的風險意識。
\end{frame}


\begin{frame}{防禦提示的效果與結論}
    \small
    \begin{itemize}
        \item 測試兩種 defense prompt：簡單版（列出 10 種失敗模式）與加強版（詳細描述加範例）
        \item 對弱模型（Qwen2.5-7B）幾乎無效；對 GPT-4o 等較強模型有提升，但有限
        \item Claude-3.5-Sonnet 即使用加強版仍低於 70\%，且增加 context 長度
    \end{itemize}

    \begin{block}{結論}
        沒有任何 agent 總分超過 60\%。僅靠修改 prompt 無法解決 agent 安全問題，作者主張未來需要更有效的方法，例如\alert{微調}。
        限制：測試案例以常識推理為主，需要專業領域知識的情境留待未來。
    \end{block}
\end{frame}


% % ############################################################
% % Summary
% % ############################################################
% \section{總結}

% \begin{frame}{三篇論文小結}
%     \tabsetup
%     \footnotesize
%     \begin{table}
%         \centering
%         \begin{tabular}{@{}l l p{6.8cm}@{}}
%             \toprule
%             論文 & 性質 & 核心重點 \\
%             \midrule
%             2510.23883 & Survey & 威脅分類 5 類；hybrid 防禦；process-aware 與 pass$^{\wedge k}$ 評估 \\
%             2601.05293 & Survey + 原型 & 資安雙面性；自主等級與有界自主；四層風險模型 \\
%             2412.14470 & Benchmark & 16 個 agent 皆 $<$ 60\%；缺乏穩健性與風險意識；defense prompt 不足 \\
%             \bottomrule
%         \end{tabular}
%     \end{table}

%     \begin{block}{共通觀察}
%         自主性愈高風險愈大；單一防禦不足，需分層組合（沙箱、政策執行、監控、人類監督）；評估需更貼近真實並看到過程層級。
%     \end{block}
% \end{frame}

\end{document}